\begin{abstract}
Housing corporations increasingly rely on image archives to document rental property conditions, yet manual inspection remains time consuming and subjective. Advances in computer vision and multi modal artificial intelligence create opportunities to improve housing inspection through automated image analysis.

Existing research often evaluates supervised object detection models and multi modal foundation models separately using different datasets, metrics, and evaluation settings. Direct comparison between these approaches in housing inspection remains limited, creating uncertainty regarding robustness, generalization, and practical suitability.

Supervised object detection approaches such as YOLOv8 demonstrate strong defect localization capabilities, while multi modal foundation models such as LLaVA 1.6 show strong visual reasoning and zero shot capabilities. Related domains including medical imaging, infrastructure monitoring, and anomaly detection already demonstrate substantial progress through artificial intelligence driven inspection methods. However, transfer of these developments toward housing inspection remains limited.

This thesis introduces a controlled comparison framework in which YOLOv8 and LLaVA 1.6 are evaluated using identical housing inspection datasets, shared inspection categories, and harmonized outputs to enable fair comparison.

The study contributes an inspection oriented benchmarking framework supporting future automated housing inspection systems and helps transfer methodological progress from mature artificial intelligence domains toward housing inspection.
\end{abstract}